Given $\delta>0$ and $k\ge3$, must there be an $m=m(\delta,k)$ such that, for all sufficiently large $N$, every $A,B\subseteq\{1,\ldots,N\}$ with $|A|\ge\delta N$ and $|B|\ge m$ has a nonconstant $k$-term arithmetic progression in $A$ with difference in $B-B$?

The assertion is false already for $k=3$, and hence for every $k\ge3$. We construct an infinite set $B$ and a set $A$ of natural density $1/50$ such that no nonconstant three-term progression in $A$ has difference in $B-B$. Passing to finite intervals then disproves the exact quantifiers in the question.

The known negative result is attributed to Furstenberg's \emph{Recurrence in ergodic theory and combinatorial number theory} at \url{https://www.erdosproblems.com/1185}. The following quadratic-phase construction supplies a direct proof. Its only distribution input, quadratic equidistribution, is justified below.
